\input{../preamble.tex}

\lecturenumber{7}
\title{Basic IPC}
\version{2.0.0}

\begin{document}

  \begin{frame}[plain, noframenumbering]
    \titlepage
  \end{frame}

\begin{slide}

\slidetitle{Interprocess Communication (IPC)}

    Processes do not share any memory with each other
    \bigskip

    Some processes might want to work together for a task, so need to communicate information
    \bigskip
    
    IPC are mechanisms to share information between processes

\end{slide}

\begin{slide}

\slidetitle{Shared Memory}

\end{slide}

\begin{slide}

    \slidetitle{Pipe}
    
    \vspace{-3mm}
    
    A pipe is a {half-duplex} communication mechanism:
    \begin{itemize}
        \item Data written into \texttt($fd[1]$) can be read from \texttt($fd[0]$).
        \item Data flows in one direction only.
    \end{itemize}
    \bigskip
    
    Regular pipes:
    \begin{itemize}
        \item Both file descriptors exist in the {same process}.
        \item The parent and child process {share the pipe} after a {fork()}.
        \item The {parent writes} to the pipe, and the {child reads} from it (or vice versa).
    \end{itemize}
    \bigskip
    
    Named pipes (FIFOs):
    \begin{itemize}
        \item A named pipe allows communication between {unrelated processes}.
        \item Exists as a {special file} in the filesystem.
    \end{itemize}

\end{slide}

\begin{slide}

    \slidetitle{Pipe Example}
    
    \begin{minted}{c}
int fd[2]; // File descriptors for the pipe
pipe(fd);

pid_t pid;
char write_msg[] = "Hello from parent!";
char read_msg[100];
pid = fork(); 

if (pid > 0) {  
    write(fd[1], write_msg, strlen(write_msg) + 1);
} else {  
    read(fd[0], read_msg, sizeof(read_msg));
    printf("Child received: %s\n", read_msg);
}
\end{minted}
    
\end{slide}

\begin{slide}

    \slidetitle{Pipe Example}
    
    \begin{minted}{c}
int fd[2]; // File descriptors for the pipe
pipe(fd);

pid_t pid;
char write_msg[] = "Hello from parent!";
char read_msg[100];
pid = fork(); 
if (pid > 0) {  
    close(fd[0]);
    write(fd[1], write_msg, strlen(write_msg) + 1);
    close(fd[1]);
} else {  
    close(fd[1]);
    read(fd[0], read_msg, sizeof(read_msg));
    close(fd[0]);
    printf("Child received: %s\n", read_msg);
}
\end{minted}
    
\end{slide}

  \begin{slide}

    \slidetitle{IPC is Transferring Bytes Between Two or More Processes}

    Reading and writing files is a form of IPC
    \medskip

    The read and write system calls allow any bytes

  \end{slide}

\begin{slide}

    \slidetitle{Signals are a Form of IPC that Interrupts}

    You could also press Ctrl+C to stop a program

    \leftspace{}This interrupts your programs execution and exits early
    \medskip

    Kernel sends a number to your program indicating the type of signal

    \leftspace{}Kernel default handlers either ignore the signal or terminate
    your process
    \medskip

    Ctrl+C sends \texttt{SIGINT} (interrupt from keyboard)
    \medskip

    If the default handler occurs the exit code will be 128 + signal number

\end{slide}

  \begin{slide}

    \slidetitle{You Can Set Your Own Signal Handlers with \texttt{sigaction}}


    You just declare a function that doesn't return a value, and has an \texttt{int} argument

    \leftspace{}The integer is the signal number
    \medskip

    Some numbers are non-standard, below are a few from Linux x86-64:
    \begin{itemize}
      \item 2: \texttt{SIGINT} (interrupt from keyboard)
      \item 9: \texttt{SIGKILL} (terminate immediately)
      \item 11: \texttt{SIGSEGV} (memory access violation)
      \item 15: \texttt{SIGTERM} (terminate)
    \end{itemize}

  \end{slide}


  \begin{slide}

    \slidetitle{Most Operations Are Non-Blocking}

    A non-blocking call returns immediately, and you check if something occurs
    \medskip

    To turn \texttt{wait} into a non-blocking call, use \texttt{waitpid} with
    \texttt{WNOHANG} in \texttt{options}
    \medskip

    To react to changes to a non-blocking call, we can either use a
    \textit{poll} or \textit{interrupt}

  \end{slide}

  \begin{slide}

    \slidetitle{Polling Continuously Calls the Function and Checks for Changes}

    We call \texttt{waitpid} over and over until the child exits

    \leftspace{}Note: some hardware behaves like this,

    \leftspace{}\leftspace{}the kernel may have to check for changes
    \bigskip

    What's the drawback of this approach?

  \end{slide}

  \begin{slide}

    \slidetitle{An Interrupt Instead Occurs Right After the Change}

    Instead of calling \texttt{wait} or \texttt{waitpid} from \texttt{main},
    we can do it in the interrupt handler

    \leftspace{}The kernel sends the \texttt{SIGCHLD} whenever one of its
                  children exit
    \medskip

    This idea also applies to the kernel, hardware can generate interrupts

  \end{slide}

  \begin{slide}

    \slidetitle{Interrupt Handlers Run to Completion}


    An interrupt may occur while an interrupt handler is already running
    \bigskip

    All interrupt handler code must be reentrant

    \leftspace{}You need to be able to pause execution,

    \leftspace{}execute another call (to the same function),

    \leftspace{}and resume execution

  \end{slide}

  \begin{slide}

    \slidetitle{On a RISC-V CPU, There's 3 Terms for ``Interrupts''}

    \textit{Interrupt}

    \leftspace{}Triggered by external hardware,

    \leftspace{}handled by the kernel (needs to respond quickly)
    \bigskip

    \textit{Exception}

    \leftspace{}Triggered by an instruction (divide by zero, illegal memory
                  access),

    \leftspace{}default handler is the kernel (calling process suspended),

    \leftspace{}the process can optionally handle some of these themselves
    \bigskip

    \textit{Trap}

    \leftspace{}Transfer of control to a trap handler caused by either

    \leftspace{}an exception or an interrupt (code that runs)
    \medskip

    A system call would be a \textit{requested trap}
  \end{slide}

\begin{slide}

	\slidetitle{Review Question}

\begin{minted}{c}
int main(){
    int fd[2];
    pid_t pid;
    char buffer[20];

    pid = fork(); 
    pipe(fd);

    if (pid == 0) {  
        write(fd[1], "Hello", 6);
        printf("Child wrote to pipe\n");
    } else {  
        read(fd[0], buffer, sizeof(buffer));
        printf("Parent received: %s\n", buffer);
    }
}
\end{minted}
	
\end{slide}

\end{document}
